% =============================================================================
\section{Motivation and the use of AI}
\label{sec:history}

The question of this paper came from quantum geometry, and I reached quantum geometry
from an unexpected direction. Since late June 2026 I had been building an attention layer
whose edges carry operators on the Albert algebra, the exceptional Jordan algebra, and whose
pair interactions are classified in advance: the tensor square of a traceless element splits
into five irreducible channels, $26\otimes26=1\oplus26\oplus52\oplus273\oplus324$
\cite{FitchPairChannels}. In October that work led me to the quantum tetrahedron. The five
channels turned out to label a basis of the shapes of a quantum tetrahedron whose four faces
carry the exceptional group $F_4$. For $\mathrm{SU}(2)$ faces the orientation operator of
that tetrahedron is the triple product $J_1\cdot(J_2\times J_3)$, the four-valent volume
operator of loop quantum gravity. The volume operator then took me to the dynamics of loop
quantum gravity, where it enters the construction of the Hamiltonian constraint. There I
tried self-dual variables, in which the edge labels belong to $\mathfrak{su}(1,1)$ instead
of $\mathfrak{su}(2)$. In a recoupling basis the triple product becomes the zero-diagonal
Jacobi matrix with cubically growing links mentioned in the introduction. It is in the
limit-circle case. So the volume operator is not determined until a boundary condition at
infinity is chosen, and different choices give different spectra.

That raised the question this paper answers in general: given a Jacobi matrix in the
limit-circle case, where exactly are the eigenvalues of each of its self-adjoint
extensions? For the cubic chains the answer first appeared numerically, as a
Bohr--Sommerfeld rule in which the extension angle enters as a phase. The rest of the paper
is what was needed to prove that rule, and to see how far it extends: to every growth
exponent $\alpha>1$, to dimerised links, and to the moment problems of \Cref{sec:moments}.
The companion paper on the volume gap~\cite{FitchVolumeGap} grew from the same operator,
with $\mathfrak{su}(2)$ labels. \Cref{app:provenance} describes how the results were
reached.

\paragraph{Use of AI.} I posed the questions, chose the problems and the order in which to
pursue them, and set the standard a result had to meet. I directed the computations and
decided what this paper claims. Large language models were used throughout. Anthropic's
Claude was used in two ways. In sessions on claude.ai it carried out the earlier
calculations on quantum tetrahedra and began the calculation on the dynamics, from which
this paper's question came. Through Claude Code it wrote and ran the programs of
\Cref{app:repro}, ran the agent campaigns, and drafted and revised the text. It also gave
the first written proofs of the results here, among them the rotation identity and the
discrete Bohr--Sommerfeld rule (\Cref{thm:discreteBS}), the discrete Olver bound
(\Cref{thm:olver}), the second-order term (\Cref{thm:second}), the dimerised law
(\Cref{thm:dimer}) and the critical indices (\Cref{thm:valent}). Separate Claude agents
refereed each result as it was written, and an independent Claude referee read the whole
paper. OpenAI's GPT models were used through ChatGPT and Codex. ChatGPT produced the first
document on the tetrahedron in the Albert algebra. A Codex session independently checked
the rate and the counting law, and another compared \Cref{cor:squeeze} with directly
computed spectra. GPT also reviewed the paper. Google's Gemini reviewed only the companion
paper and had no part in this one. Review comments were adopted only after they had been
checked against the source and our own computations. The claims rest on the following, and
not on any model's assertion: the written proofs in this paper, none of which depends on a
computation; exact symbolic checks of the identities listed in \Cref{app:repro}; the
independent eigenvalue computations of \Cref{sec:numerics}; and cited results.
Responsibility for the content, and for any error that remains, is mine.
